\section{Phase 1D: RLAIF Reward and Preference Pipeline}

Phase 1D converts BudgetRAG logs into normalized RLAIF-style datasets:

\begin{center}
\small
\begin{tabular}{>{\raggedright\arraybackslash}p{4.7cm} >{\raggedright\arraybackslash}p{8.5cm}}
\toprule
Command & Role \\
\midrule
\code{rlaif-build} & Convert BudgetRAG rows into normalized actions and feedback. \\
\code{rlaif-label-answers} & Label answer correctness, support, faithfulness, and unsupported claims. \\
\code{rlaif-label-contexts} & Label context sufficiency, useful chunks, redundant chunks, irrelevant chunks, and missing evidence. \\
\code{rlaif-reward} & Build scalar rewards and pairwise preferences with quality guardrails. \\
\code{rlaif-split} & Split by \code{benchmark + query\_id} to avoid row-level leakage. \\
\code{rlaif-train/eval} & Train and evaluate offline selector baselines. \\
\code{rlaif-label-pairs} & Directly judge action pairs for reward calibration audits. \\
\bottomrule
\end{tabular}
\end{center}

The reward builder preserves provenance and missing reasons. Missing, ambiguous, or invalid labels do not become score zero. Preferences are generated within comparable query groups, with guardrails that reject efficiency-only wins when answer quality drops too much.

RLAIF in this report means AI-feedback labeling for reward/preference construction. It is not full RLHF and it is not human preference learning.
